% Chapter 18 -- Claim: Low-Risk Stocks Beat High-Risk Stocks
% Requires: amsmath, amssymb, booktabs, array, tabularx, longtable, graphicx, tcolorbox (breakable, skins), hyperref or url.
% Defines the callout box only if the preamble has not.
\providecolor{fenavy}{HTML}{1B2A6B}
\providecolor{feshade}{HTML}{F3F5FA}
\makeatletter
\@ifundefined{callout}{%
  \newtcolorbox{callout}{breakable, enhanced,
    colback=white, colframe=gray!70, boxrule=0.5pt, arc=1mm,
    borderline west={2.5pt}{0pt}{gray!70},
    left=8pt, right=8pt, top=4pt, bottom=4pt}%
}{}
\makeatother

\providecolor{fenavy}{HTML}{1B2A6B}
\providecolor{feshade}{HTML}{F3F5FA}
\makeatletter
% Shaded box for optional technical material that general readers may skip.
\@ifundefined{technote}{%
  \newtcolorbox{technote}{breakable, enhanced,
    colback=feshade, colframe=fenavy!60, boxrule=0.5pt, arc=1mm,
    left=8pt, right=8pt, top=4pt, bottom=4pt}%
}{}
\makeatother

\chapter{Claim: Low-Risk Stocks Beat High-Risk Stocks}
\label{ch:15}
\providecommand{\term}[1]{\emph{#1}}

Common sense says that riskier investments must pay more on average, or no one would hold them. Yet investors who bought the dull, steady stocks of the market, such as utilities and grocery chains, have historically done about as well as those who bought the most exciting ones, with far smaller swings along the way. This chapter asks whether that pattern is real, why it might exist, and whether an investor can profit from it.

The capital asset pricing model\index{capital asset pricing model (CAPM)} makes one central prediction: assets that move more with the market should earn more. If a stock's beta\index{beta} doubles, its expected excess return\index{return!excess} should double too. This was among the first predictions in finance to be tested, and Chapter~\ref{ch:03} showed that it failed early. Black, Jensen and Scholes (1972) and Fama and MacBeth (1973), who estimated the link between returns and beta month by month and averaged the results, found that average returns did rise with beta, but far less steeply than the theory implied. Black (1972) offered an explanation based on limits to borrowing. Some investors want more risk but cannot borrow freely. Instead of borrowing, they buy high-beta stocks, which pushes up their prices and lowers their later returns.

For two decades the flat security market line was treated as a curiosity. It became a claim about profits when later research made the effect larger and easier to trade. Ang, Hodrick, Xing and Zhang (2006) tested the theory that only market risk is priced, and found that stocks with high \term{idiosyncratic volatility}\index{volatility}, the part of volatility not explained by the market, earned strikingly low returns. Asking why professionals did not trade such mispricing away, Baker, Bradley and Wurgler (2011) argued that fund managers judged against a benchmark have little reason to correct the mispricing. To show that the pattern reflects limits on borrowing everywhere, not a quirk of U.S. stocks, Frazzini and Pedersen (2014) built a leveraged ``betting-against-beta'' factor that earned high risk-adjusted returns in many countries and asset classes. The pattern is now called the \term{low-risk anomaly}\index{low-risk anomaly|textbf}. Novy-Marx and Velikov (2022) later argued that much of the betting-against-beta performance comes from the way the factor is built, not from low risk itself.

The question matters in practice. Low-volatility funds grew rapidly after 2010, and companies still use beta to set the return a new project must earn. The anomaly is also hard to exploit, because exploiting it needs leverage, which is exactly what the explanation says investors avoid. The claim on trial\index{claims-on-trial method} is: \textbf{low-beta and low-volatility stocks earn higher risk-adjusted returns than high-risk stocks, and a leveraged long--short portfolio captures the difference.}

The lab in Section~\ref{sec:18-lab} repeats two classic tests on the free beta- and variance-sorted portfolios of the Kenneth R. French Data Library: the security market line of Black, Jensen and Scholes (1972) and the cross-sectional regressions of Fama and MacBeth (1973). It then rebuilds the betting-against-beta factor of Frazzini and Pedersen (2014) and compares it with AQR's free published series. The goal is to check whether low-risk stocks still earn more than the CAPM predicts once costs and construction choices are taken into account.

By the end of the chapter you will be able to:

\begin{itemize}
  \item estimate a beta, split total risk\index{risk!total} into systematic and idiosyncratic parts, and compute a CAPM alpha\index{alpha};
  \item draw the security market line\index{security market line}, and explain how limits on leverage\index{leverage} flatten it;
  \item build a beta-neutral\index{market-neutral} betting-against-beta\index{betting against beta (BAB)} portfolio, and compare the explanations for the anomaly, such as leverage limits and lottery preferences\index{lottery preference};
  \item judge whether a minimum-volatility or BAB fund can deliver what its backtest\index{backtest} promises.
\end{itemize}


\section{Foundations: beta, the CAPM and the security market line}
\label{sec:18-foundations}

\subsection{Beta, the CAPM and the security market line}

Chapter~\ref{ch:03} introduced beta, the CAPM and the security market line (Section~\ref{sec:04-capm}). Beta is the slope of a regression of an asset's excess return on the market's excess return (Equation~\eqref{eq:04-capm}). The market itself has a beta of 1. The CAPM says that expected excess returns rise in a straight line with beta. Two consequences of the model matter for this chapter.

First, the security market line passes through the risk-free rate\index{risk-free rate} at $\beta = 0$ and through the market at $\beta = 1$, so its slope equals the market premium.

Second, an asset's \term{total risk}\index{risk!total|textbf}, the variance\index{variance} of its return, splits into two parts. The systematic part is $\beta_i^2\,\mathrm{Var}(R_m)$. The idiosyncratic part is the variance of the regression residual\index{regression!residuals}. Only the systematic part earns a reward, because investors can diversify\index{diversification} the idiosyncratic part away at no cost. Two stocks with the same volatility can therefore have very different betas, and the CAPM says they should have very different expected returns\index{return!expected}.

\subsection{Alpha}

An asset's alpha, often called Jensen's alpha, is the intercept of the same regression: the average return beyond what beta explains (Section~\ref{sec:04-capm}). Under the CAPM every alpha is zero. A multifactor alpha adds other factors to the regression, such as size, value and momentum\index{momentum} (Chapter~\ref{ch:03}). The low-risk anomaly is a claim about how alpha varies with beta: low-beta stocks have had positive alphas and high-beta stocks negative ones.

The CAPM assumes that investors can borrow and lend freely at the risk-free rate. An investor who wants more risk than the market can then borrow and buy more of the market. Many investors cannot do this. Mutual funds and pension plans face limits on \textbf{leverage}\index{leverage|textbf}, which means borrowing to hold more than their capital. Individuals face \textbf{margin} requirements and pay more than the risk-free rate to borrow. Black (1972) showed that without riskless borrowing the SML is still a straight line, but a flatter one. It starts at the return on a \term{zero-beta portfolio}\index{zero-beta portfolio|textbf}, a portfolio uncorrelated with the market, whose return lies above the risk-free rate.

\textbf{Worked example.} The risk-free rate is 3\% and the market premium is 6\% (Figure~\ref{fig:18-sml}).

\begin{itemize}
  \item Stock A has a beta of 0.6 and has averaged 9\%. The CAPM predicts $3\% + 0.6 \times 6\% = 6.6\%$, so A's alpha is $+2.4\%$.
  \item Stock B has a beta of 1.5 and has averaged 10\%. The CAPM predicts $3\% + 1.5 \times 6\% = 12\%$, so B's alpha is $-2\%$.
  \item An investor who can borrow at 3\% can lever A to a beta of 1. She holds $1/0.6 = 1.67$ dollars of A for each dollar of capital and borrows the remaining 0.67 dollars. Her expected return is $3\% + 1.67 \times (9\% - 3\%) = 13\%$, compared with 9\% for the market at the same beta. Even at a borrowing rate of 5\%, she earns $1.67 \times 9\% - 0.67 \times 5\% \approx 11.7\%$.
  \item An investor who cannot borrow and wants a beta above 1 must buy B instead. If many investors are in that position, they bid up high-beta stocks such as B and push down their future returns. That is the core of the leverage-constraint explanation in Section~\ref{sec:18-explanations}.
\end{itemize}

\begin{figure}[htbp]
\centering
\scalebox{0.9}{%
\begin{tikzpicture}[
  % CAPM: r = 3 + 6 * beta. A = (0.6, 9), B = (1.5, 10); the line through A and B
  declare function={capm(\b) = 3 + 6*\b;  flat(\b) = 9 + (\b - 0.6)/0.9;}
]
\begin{axis}[
  width=0.6\linewidth, height=0.34\linewidth,
  axis lines=left, axis line style={-stealth},
  xmin=0, xmax=2, ymin=0, ymax=15.5,
  xtick={0,0.5,...,2}, ytick={0,5,10,15},
  xticklabel style={/pgf/number format/fixed},
  xlabel={Beta}, ylabel={Average return (\%)},
  clip=false,
]
  \addplot[fegray, thick, dashed, domain=0:2, samples=2] {flat(x)};
  \addplot[fenavy, very thick, domain=0:2, samples=2] {capm(x)};
  \fill (axis cs:0.6,9)  circle (2.6pt);
  \fill (axis cs:1,9)    circle (2.6pt);
  \fill (axis cs:1.5,10) circle (2.6pt);
  \node[anchor=south] at (axis cs:0.6,9.5) {A: $\alpha = +2.4\%$};
  \node[anchor=north west] at (axis cs:0.98,8.4) {market};
  \node[anchor=north west] at (axis cs:1.47,9.1) {B: $\alpha = -2\%$};
  \node[anchor=north west] at (axis cs:0.02,2.7) {$r_f = 3\%$};
  \node[fenavy, anchor=west] at (axis cs:2.02,15) {CAPM line};
  \node[fegray, anchor=west, align=left] at (axis cs:2.02,{flat(2)})
    {line through\\A and B};
\end{axis}
\end{tikzpicture}%
}
\caption{Low-beta stocks have done better, and high-beta stocks worse, than the CAPM predicts. The security market line in the worked example. The CAPM line runs from the risk-free rate of 3\% through the market at a beta of 1. Stock A lies above it (positive alpha) and stock B below it (negative alpha), so the line through A and B is flatter than the CAPM predicts.}
\label{fig:18-sml}
\end{figure}

\section{How it works}
\label{sec:18-how}

There are three ways to sort stocks on risk, and the literature uses all three.

\begin{itemize}
  \item \textbf{Beta.}\index{beta|textbf} Estimate each stock's beta from past returns, typically 36 to 60 months of monthly data or a year of daily data. Sort into portfolios. This tests the CAPM directly.
  \item \textbf{Total volatility.} Sort on the standard deviation\index{standard deviation} of past daily returns, often over one to three months. This needs no market model, and it is how most low-volatility funds work.
  \item \textbf{Idiosyncratic volatility.}\index{idiosyncratic volatility|textbf} Sort on the volatility of residuals from a factor model\index{factor model}, usually the Fama--French\index{Fama--French factors} three-factor model. Under the CAPM this risk should not be priced at all, so a negative relation is especially puzzling.
\end{itemize}

The three measures are strongly correlated. High-beta stocks tend to be volatile, and volatile stocks tend to have high idiosyncratic risk\index{risk!idiosyncratic}. Separating them is one of the open problems in this chapter.

A flat SML shows up in two ways. The first is in \emph{alphas}: low-risk portfolios have positive CAPM alphas, and high-risk portfolios negative ones. The second is in a \emph{long-short portfolio\index{long--short portfolio}}. The simple version buys low-risk stocks and sells high-risk stocks short. It has a strongly negative beta, so it loses money in rising markets even when the SML is flat. Its average return can be negative while its alpha is positive. \term{Betting against beta (BAB)}\index{betting against beta (BAB)|textbf} fixes this. It levers the low-beta leg up to a beta of 1 and scales the high-beta leg down to a beta of 1. The portfolio is then \term{beta-neutral}\index{market-neutral|textbf}, with no net exposure to the market (Figure~\ref{fig:18-bab}). The shaded box gives the formula.

\begin{technote}
\textbf{Technical note: the BAB return.} In symbols, the portfolio's excess return is
\begin{equation}
  r_{\mathrm{BAB}} = \frac{r_L - r_f}{\beta_L} - \frac{r_H - r_f}{\beta_H}
  \label{eq:18-bab}
\end{equation}
If the SML is flat, this portfolio earns a positive return with no market exposure. The betas must be \term{ex ante}, estimated from data available when the portfolio is formed. Otherwise the construction looks ahead (Section~\ref{sec:02-lookahead}).
\end{technote}

\begin{figure}[htbp]
\centering
\scalebox{0.9}{%
\begin{tikzpicture}
\begin{axis}[
  width=0.6\linewidth, height=0.34\linewidth,
  ybar, bar width=0.6cm,
  axis lines=left, axis line style={-stealth},
  ymin=-2, ymax=2.2, ytick={-2,-1,0,1,2},
  ylabel={Dollars per \$1 of capital},
  symbolic x coords={L,H}, xtick=data,
  xticklabels={{Low-beta leg\\($\beta_L = 0.6$)}, {High-beta leg\\($\beta_H = 1.5$)}},
  xticklabel style={align=center},
  enlarge x limits=0.45,
  nodes near coords,
  nodes near coords style={font=\small},
  extra y ticks={0}, extra y tick labels={},
  extra y tick style={grid=major, major grid style={gray, thin}},
  legend style={draw=none, at={(0.5,1.03)}, anchor=south, legend columns=2,
                /tikz/every even column/.append style={column sep=1em}},
]
  \addplot[fill=fenavy, draw=none] coordinates {(L,1.67) (H,-0.67)};
  \addplot[fill=gray!40, draw=none] coordinates {(L,1) (H,-1)};
  \legend{position size, resulting beta exposure}
\end{axis}
\end{tikzpicture}%
}
\caption{Betting against beta borrows to scale up low-beta stocks and scales down high-beta ones until the market's moves cancel. Building a BAB portfolio from the worked example. Per dollar of capital, it holds \$1.67 of the low-beta stock ($\beta_L = 0.6$) and shorts \$0.67 of the high-beta stock ($\beta_H = 1.5$). Each leg then carries a beta exposure of 1, in opposite directions, so the portfolio's beta is zero.}
\label{fig:18-bab}
\end{figure}

The explanations fall into two groups. One treats the anomaly\index{anomaly} as \term{mispricing}\index{mispricing|textbf}. The other treats it as \textbf{compensation for a risk the CAPM misses}. Section~\ref{sec:18-explanations} sets them out. The difference matters because the two groups predict different things: when the anomaly should be strongest, and whether it should survive once it becomes popular.

\section{The claim}
\label{sec:18-claim}

\begin{callout}
\textbf{Claim.} Low-beta and low-volatility stocks earn higher risk-adjusted returns than the CAPM predicts. The security market line is flatter than the market premium, and a beta-neutral portfolio that is long low-risk stocks and short high-risk stocks earns a positive alpha.
\end{callout}

To test it we need to specify four things.

\begin{itemize}
  \item \textbf{The strong form.}\index{efficient market hypothesis!strong form|textbf} The SML is \emph{inverted\index{yield curve!inverted}}: low-risk stocks earn higher raw returns. Some volatility sorts find this; most beta sorts find a flat or weakly positive relation.
  \item \textbf{The weak form.}\index{efficient market hypothesis!weak form|textbf} The SML is positive but flatter than the CAPM predicts, so low-beta stocks have positive alphas, as Black, Jensen and Scholes found. This is the claim tested here.
  \item \textbf{Risk adjustment.} A CAPM alpha is the first test; the claim is stronger if the alpha survives size, value and momentum, and stronger still if it survives profitability and interest-rate exposure.
  \item \textbf{What would count against it.} An alpha concentrated in micro-caps\index{micro-caps}, explained by other factors or faded since the strategy became popular. The claim fails if a tradable beta-neutral version earns nothing after costs.
\end{itemize}

\section{The original evidence}
\label{sec:18-evidence}

\textbf{A flat line from the start.} Because individual betas are measured with large errors\index{estimation error}, Black, Jensen and Scholes (1972) sorted NYSE stocks from 1931 to 1965 into ten portfolios by past beta. Average returns rose with beta, but much less than the CAPM predicted, so low-beta portfolios had positive alphas. Fama and MacBeth (1973) confirmed a positive but modest slope, and Haugen and Heins (1975) found that low-volatility portfolios had matched or beaten high-volatility ones. Because size and beta are closely related, Fama and French (1992) included size and found almost no relation between beta and average return from 1963 to 1990.

\textbf{Idiosyncratic volatility.}\index{idiosyncratic volatility|textbf} Ang, Hodrick, Xing and Zhang (2006) found that the fifth of U.S. stocks with the highest idiosyncratic volatility earned 1.06\% a month less than the lowest fifth, a strong form of the anomaly, and later found the same in 23 developed markets.

\textbf{From anomaly to strategy.} Baker, Bradley and Wurgler (2011) argued that \term{benchmarking} stops institutions from exploiting the anomaly. Frazzini and Pedersen (2014) built the BAB factor of Section~\ref{sec:18-how}, which had a Sharpe ratio\index{Sharpe ratio} of 0.78 in U.S. stocks from 1926 to 2012, alphas that survived size, value, momentum and liquidity\index{liquidity} controls, and success in 19 international markets and in bonds and futures. It did worse when borrowing was hard, as their leverage theory predicts. Low-volatility index funds launched around 2011 and grew quickly.

\section{Explanations}
\label{sec:18-explanations}

\textbf{Leverage constraints.}\index{leverage|textbf} Investors who want high returns but cannot borrow buy high-beta stocks instead of a levered market portfolio, which bids them up and flattens the SML (Black 1972; Frazzini and Pedersen 2014). As predicted, constrained investors hold high-beta portfolios, and BAB does worse when borrowing becomes harder.

\textbf{Lottery preferences.}\index{lottery preference|textbf} Some investors overpay for stocks with a small chance of a very large gain. Barberis and Huang (2008) derive this from \term{prospect theory}\index{prospect theory}, in which people overweight unlikely outcomes, and Bali, Cakici and Whitelaw (2011) found that stocks with the largest recent one-day returns underperform, which explains much of the idiosyncratic volatility effect.

\textbf{Benchmarking.} For a manager judged against a market-weighted index, a low-beta stock with a positive alpha adds \term{tracking error}\index{tracking error|textbf}, the volatility of the gap between the fund's return and the benchmark's, and lags in rising markets, so institutions do not correct the mispricing (Baker, Bradley and Wurgler 2011).

\textbf{Overpricing that is hard to short.} Liu, Stambaugh\index{Stambaugh bias} and Yuan (2018) found that the beta anomaly comes from beta's correlation\index{correlation} with idiosyncratic volatility and is concentrated among overpriced stocks, where short-selling\index{short selling} is costly.

\textbf{Missing risk factors.} Low-risk stocks tend to be large, profitable and conservative, and the five-factor model explains much of their alpha (Fama and French 2016). They also behave partly like bonds, and an interest-rate premium explains part of the anomaly, flattered by falling yields from 1981 to 2020 (Driessen, Kuiper, Nazliben and Beilo 2019). And modelling time-varying betas\index{beta!time-varying} largely removes the alpha gap (Cederburg and O'Doherty 2016).

These explanations can all be partly true. Asness, Frazzini, Gormsen and Pedersen (2020) split BAB into a bet against correlation, linked to \term{leverage constraints}, and a bet against volatility, linked to lottery demand, and found that both earn returns.

\section{Lab~\thechapter: Replicating Black, Jensen and Scholes (1972) and Frazzini and Pedersen (2014)}
\label{sec:18-lab}

\textbf{Claim on trial\index{claims-on-trial method}:} low-risk stocks earn more than the CAPM predicts. The lab replicates Black, Jensen and Scholes (1972), Fama and MacBeth (1973) and Frazzini and Pedersen (2014), using ten portfolios formed on market beta, on total variance and on residual variance, each value- and equal-weighted\index{index weighting!equal-weighted}, from July 1963, from the Kenneth French\index{Kenneth French data library} library, with the Fama--French and momentum factors and the 10-year Treasury yield from FRED\index{FRED}. Unlike Frazzini and Pedersen, who use individual stocks with rank weights, the lab uses published portfolios that anyone can reproduce, so expect the same pattern with smaller numbers.

\textbf{Part A, the SML.} For each set of deciles, average excess return, realized beta, CAPM and four-factor alphas, and the slope of return on beta, plotted against the CAPM line (Lab Figure~\ref{ch:15}.1).

\textbf{Part B, long-short portfolios.} A low-minus-high portfolio, long the three lowest-risk deciles and short the three highest, and a BAB-style portfolio that levers each leg to a beta of one, with betas estimated from the previous 60 months (Lab Figure~\ref{ch:15}.2).

\textbf{Part C, robustness checks.} Equal versus value weights, three subperiods, interest-rate exposure, returns in 2008, 2020 and 2022, and rolling ten-year alphas (Lab Figures~\ref{ch:15}.3 and~\ref{ch:15}.4).

\textbf{Part D, the investor's view.} The lowest-volatility and lowest-beta quintiles against the market: Sharpe ratio, alpha, drawdown\index{drawdown}, tracking error and the share of years they trailed (Lab Figure~\ref{ch:15}.5).

\label{sec:18-stress}
\textbf{Robustness checks.} Novy-Marx and Velikov (2022) found that AQR's BAB factor puts on average about \$1.05 of each dollar into the smallest 1\% of stocks\index{micro-caps|textbf}, which are expensive to trade (Chapter~\ref{ch:14}), so if the lab's effect is much stronger with equal than value weights\index{index weighting!value-weighted}, small stocks are doing much of the work (Section~\ref{sec:02-microcaps}). The lab uses four factors, so its alphas probably overstate what the five-factor model leaves. 2022, when stocks and bonds fell together, is the natural out-of-sample\index{in-sample versus out-of-sample} test of interest-rate exposure, and the post-2010 subperiod shows whether crowding\index{crowding|textbf} has shrunk the premium. Finally, stocks sorted into low past-beta portfolios tend to have higher true betas, which flattens the SML even if the CAPM holds\index{estimation error|textbf}. The lab ends by asking you to reach your own verdict on the scale of Section~\ref{sec:01-method} and compare it with Section~\ref{sec:18-verdict}.

\section{Where the evidence stands}
\label{sec:18-verdict}

The published studies point to a clear pattern: \textbf{the security market line is flatter than the CAPM predicts, one of the most robust findings in empirical finance. Whether low risk earns a separate premium that a beta-neutral portfolio can collect depends on how the portfolio is built, on small stocks and on overlap with other factors.} The table summarizes the evidence as of November 2026.

{\small
\begin{longtable}{@{}>{\raggedright\arraybackslash}p{0.19\linewidth}>{\raggedright\arraybackslash}p{0.13\linewidth}>{\raggedright\arraybackslash}p{\dimexpr 0.45\linewidth-6\tabcolsep\relax}>{\raggedright\arraybackslash}p{0.23\linewidth}@{}}
\toprule
Situation & Holds? & Why & Evidence \\
\midrule
\endfirsthead
\multicolumn{4}{@{}l}{\emph{(continued)}} \\
\toprule
Situation & Holds? & Why & Evidence \\
\midrule
\endhead
\midrule
\multicolumn{4}{r@{}}{\emph{(continued on the next page)}} \\
\endfoot
\bottomrule
\endlastfoot
A flat security market line & \textbf{Yes} & Returns rise with beta far less than the market premium implies, in the U.S. and abroad. & Black, Jensen and Scholes (1972); Fama and MacBeth (1973); Fama and French (1992) \\*
\addlinespace[3pt]
Idiosyncratic volatility & \textbf{Yes} & The most volatile stocks earned about 1\% a month less than the least, in 23 markets; lottery demand absorbs much of it. & Ang, Hodrick, Xing and Zhang (2006); Bali, Cakici and Whitelaw (2011) \\
\addlinespace[3pt]
The beta-neutral BAB factor & \textbf{Partly} & High risk-adjusted returns across countries and assets, but much capital in the smallest stocks. & Frazzini and Pedersen (2014); Novy-Marx and Velikov (2022) \\
\addlinespace[3pt]
After profitability and investment factors & \textbf{Weakly} & The five-factor model explains much of the alphas. & Fama and French (2016) \\
\addlinespace[3pt]
With betas that vary over time & \textbf{Weakly} & Time-varying betas largely remove the alpha gap. & Cederburg and O'Doherty (2016) \\
\addlinespace[3pt]
Interest-rate exposure & \textbf{Partly} & An interest-rate premium explains part, but not all, of the anomaly. & Driessen, Kuiper, Nazliben and Beilo (2019) \\*
\addlinespace[3pt]
Why the line is flat & \textbf{Disputed} & Leverage constraints, lottery preferences, benchmarking and costly short selling all have support. & Frazzini and Pedersen (2014); Baker, Bradley and Wurgler (2011); Liu, Stambaugh and Yuan (2018) \\
\end{longtable}
\addtocounter{table}{-1}% uncaptioned longtable: keep table numbers unchanged
}

Low-volatility portfolios have delivered returns close to the market's with less volatility, which is useful to an unconstrained investor, but inflows since 2011 may have raised their prices and made past returns a poor guide.

\textbf{What would change the picture.} The strong form would gain support if a beta-neutral, value-weighted version kept a significant alpha against the five-factor model and interest-rate exposure out of sample. It would lose support if the post-2010 alphas proved to be inflow effects that reversed.

%\label{sec:18-open}
%\textbf{Open questions.} Which risk measure is priced, beta, as leverage theories suggest, or volatility, as lottery theories suggest? And how large is the alpha of a value-weighted, beta-neutral portfolio of liquid stocks against the five-factor model?

\section{What practitioners and investors can use}
\label{sec:18-investors}

\textbf{Minimum-volatility funds}\index{minimum-volatility funds|textbf}, long-only portfolios chosen for low volatility, are the most common way to access the anomaly. With a beta around 0.7 to 0.8, they usually trail strong bull markets and beat the market in crashes, so the case for them is a better return per unit of risk, not higher returns. Their tracking error\index{tracking error|textbf} is large enough that they can lag the market for years, and they tend to overweight interest-sensitive sectors such as utilities and consumer staples. \textbf{BAB-style funds}\index{betting against beta (BAB)|textbf} are a different product: \term{market-neutral}, levered, exposed to the micro-cap issues above, and prone to sharp losses when high-beta stocks rebound.

\begin{callout}
\textbf{What Failed: the CAPM as a hurdle rate.} The CAPM's empirical problems have been known since the early 1970s, yet it remains the standard way to set the cost of equity: Graham and Harvey (2001) found that 73.5\% of 392 U.S. chief financial officers always or almost always used it. A flat SML means the CAPM overstates the cost of capital for high-beta projects and understates it for low-beta ones, so firms may reject good low-beta projects and accept poor high-beta ones. The failure is treating a model of \emph{risk} as a reliable model of \emph{expected return} when the evidence has long said otherwise.
\end{callout}

% Key terms of the chapter, in alphabetical order, for readers new to finance.
% A term defined in an earlier chapter points to the chapter, section and page.
\providecolor{FEMaroon}{RGB}{128,30,45}
\providecolor{FECream}{RGB}{253,248,238}
\providecolor{FEGold}{RGB}{196,150,60}
\makeatletter
\@ifundefined{keyterms}{%
  \newtcolorbox{keyterms}[1][Key terms in this chapter]{breakable, enhanced,
    colback=FECream, colframe=FEGold, boxrule=0.6pt, arc=2pt, leftrule=3pt,
    left=6pt, right=6pt, top=3pt, bottom=3pt, before skip=8pt, after skip=8pt,
    fonttitle=\bfseries, coltitle=white, colbacktitle=FEMaroon, title={#1}}%
}{}
\makeatother

\begin{keyterms}
\footnotesize
\begin{description}[leftmargin=0pt, itemsep=0pt, parsep=0pt, topsep=0pt, font=\normalfont\bfseries\color{FEMaroon}]
  \item[Benchmarking.] Benchmarking means judging a fund manager against an index, which discourages holding assets that make the fund's return differ from the index.
  \item[Beta-neutral.] A portfolio is beta-neutral when its long and short sides have the same beta, so it has no net exposure to the market.
  \item[Betting against beta (BAB).] Betting against beta buys low-beta stocks levered up to a beta of 1 and sells high-beta stocks scaled down to a beta of 1.
  \item[Ex ante.] Ex ante means estimated beforehand, using only data available when the portfolio is formed.
  \item[Idiosyncratic volatility.] Idiosyncratic volatility is the part of a stock's volatility that the market does not explain.
  \item[Leverage constraints.] Leverage constraints are limits on how much an investor may borrow, which push investors who want more return toward risky stocks.
  \item[Low-risk anomaly.] The low-risk anomaly is the finding that low-beta and low-volatility stocks have earned higher risk-adjusted returns than high-risk stocks.
  \item[Market-neutral.] A market-neutral strategy is built to have no net exposure to the stock market, so its returns should not depend on whether the market rises or falls.
  \item[Mispricing.] Mispricing is a gap between a security's price and its fair value that investor errors or limits on trading allow to persist.
  \item[Prospect theory.] See Chapter~\ref{ch:13}, Section~\ref{sec:hp-prospect}, page~\pageref{kt:prospect-theory}.
  \item[Total risk.] Total risk is the whole variance of an asset's return, the sum of its systematic and idiosyncratic parts.
  \item[Tracking error.] Tracking error is the volatility of the gap between a fund's return and its benchmark's return.
  \item[Zero-beta portfolio.] A zero-beta portfolio is a portfolio uncorrelated with the market.
\end{description}
\end{keyterms}
